% !TEX program = lualatex
\documentclass[12pt,a5paper]{article}

\usepackage{polyglossia}
\setmainlanguage{french}
\usepackage{prepa}

\usepackage{geometry}
\geometry{%
  left=1cm,
  right=1cm,
  top=1cm,
  bottom=1cm
}

\title{Résolution de problème : conception d'un Thérémine}
\date{}

\begin{document}
\maketitle

Un thérémine est un instrument de musique électronique produisant un son dont la fréquence dépend de la distance entre l'instrument et la main du joueur.

Les thérémines analogiques sont constitués de deux oscillateurs identiques. La capacité d'un des oscillateurs est reliée à l'antenne et voit donc sa valeur changer faiblement (variation de l'ordre de \SI{10}{pF}). Les signaux des deux oscillateurs sont combinés pour produire un signal sonore dont la fréquence change entre $0$ et $\SI{2}{kHz}$ en fonction de la position de la main.

\textbf{Proposer un montage permettant de réaliser un Thérémine comportant des ALI, résistors, condensateurs et multiplieur.} Les valeurs des résistances et capacités seront portées sur le schéma du montage et seront choisies parmi les valeurs réalistes utilisables en TP.

Le travail se fera par groupe comportant entre 2 et 4 membres et sera d'abord réalisé au brouillon avant d'être mis au propre sur une feuille A3.

\end{document}
